\documentclass[11pt]{article}
\usepackage[utf8]{inputenc}
\usepackage[spanish]{babel}
\usepackage{amsmath, amssymb}
\usepackage{microtype}
\usepackage{geometry}
\geometry{
  top=1.8cm, bottom=2.2cm, left=2.0cm, right=2.0cm
}

% Estilo tipográfico y espaciados más compactos
\setlength{\parskip}{4pt}
\setlength{\parindent}{0pt}
\linespread{1.03}

% Figuras y captions compactos
\usepackage[font=small,labelfont=bf]{caption}
\usepackage{subcaption}
\setlength{\abovecaptionskip}{6pt}
\setlength{\belowcaptionskip}{-2pt}

% Gráficos
\usepackage{xcolor}
\usepackage{tikz}
\usepackage{pgfplots}
\pgfplotsset{compat=1.18}
\usetikzlibrary{fillbetween}
\usepackage{float}
\usepackage{hyperref}

\begin{document}

\begin{center}
{\large \textbf{Guía 6: Econometría ICS-294}}
\end{center}

\textbf{Ejercicio 1} Suponga que estima un modelo de regresión y obtiene 
$\hat{\beta}_1 = 0.56$ y valor-$p = 0.086$ en la prueba de 
$H_0: \beta_1 = 0$ contra $H_1: \beta_1 \neq 0$. 
¿Cuál es el valor-$p$ en la prueba de $H_0: \beta_1 = 0$ contra $H_1: \beta_1 > 0$?

{\color{blue}

\[
\hat{\beta}_1 = 0.56, \qquad p_{\text{bilateral}} = 0.086
\]

\textbf{Prueba bilateral: } $H_0:\beta_1=0 \ \text{vs}\ H_1:\beta_1\neq 0$.

\textbf{1. Relación entre las pruebas bilateral y unilateral}
\[
\text{valor-p}_{\text{bilateral}} \,=\, 2\,\Pr_{H_0}\!\big(T > |t_{\text{obs}}|\big), 
\qquad
\text{valor-p}_{\text{cola derecha}}=\, \Pr_{H_0}\!\big(T > t_{\text{obs}}\big),
\]
por lo tanto, si $t_{\text{obs}} > 0$, se cumple que 
$\text{valor-p}_{\text{bilateral}} = 2\,\text{valor-p}_{\text{cola derecha}}$.

\textbf{2. Cálculo del valor-$p$ unilateral}
\[
\text{valor-p}_{\text{unilateral}} \,=\, \frac{\text{valor-p}_{\text{bilateral}}}{2}
\,=\, \frac{0.086}{2} \,=\, 0.043.
\]}

\begin{figure}[H]
\centering
\begin{subfigure}{0.48\linewidth}
\centering
\begin{tikzpicture}
\begin{axis}[
    width=\linewidth,
    height=5.2cm,
    domain=-5:5,
    samples=300,
    xlabel={\small $t$},
    ylabel={\small densidad},
    ymin=0, ymax=0.42,
    title={\small Bilateral: $p=0.086$ (dos colas)},
    legend style={draw=none, fill=none, at={(0.97,0.97)}, anchor=north east, font=\footnotesize},
    ticklabel style={/pgf/number format/fixed, font=\footnotesize},
    label style={font=\footnotesize}
]
% Curva (ilustrativa con N(0,1))
\addplot[name path=curve, very thick] {1/sqrt(2*pi)*exp(-x^2/2)};
\def\tcrit{1.72}
% Colas sombreadas
\addplot[fill=blue!30, draw=none, domain=\tcrit:5] {1/sqrt(2*pi)*exp(-x^2/2)} \closedcycle;
\addplot[fill=blue!30, draw=none, domain=-5:-\tcrit] {1/sqrt(2*pi)*exp(-x^2/2)} \closedcycle;
% Líneas verticales
\addplot[dashed] coordinates {(\tcrit,0) (\tcrit,0.41)};
\addplot[dashed] coordinates {(-\tcrit,0) (-\tcrit,0.41)};
\legend{\footnotesize densidad $t$ }
\end{axis}
\end{tikzpicture}
\caption{\footnotesize Prueba bilateral $H_1\neq 0$}
\end{subfigure}\hfill
\begin{subfigure}{0.48\linewidth}
\centering
\begin{tikzpicture}
\begin{axis}[
    width=\linewidth,
    height=5.2cm,
    domain=-5:5,
    samples=300,
    xlabel={\small $t$},
    ylabel={\small densidad},
    ymin=0, ymax=0.42,
    title={\small Unilateral: $p=0.043$ (cola derecha)},
    legend style={draw=none, fill=none, at={(0.97,0.97)}, anchor=north east, font=\footnotesize},
    ticklabel style={/pgf/number format/fixed, font=\footnotesize},
    label style={font=\footnotesize}
]
\addplot[name path=curve2, very thick] {1/sqrt(2*pi)*exp(-x^2/2)};
\def\tcrit{1.72}
\addplot[fill=blue!30, draw=none, domain=\tcrit:5] {1/sqrt(2*pi)*exp(-x^2/2)} \closedcycle;
\addplot[dashed] coordinates {(\tcrit,0) (\tcrit,0.41)};
\legend{\footnotesize densidad $t$ }
\end{axis}
\end{tikzpicture}
\caption{\footnotesize Prueba unilateral $H_1:\beta_1>0$}
\end{subfigure}

\end{figure}

{\color{blue}

{\bf 3. Caso propuesto en clase }

\begin{figure}[H]
\centering
\begin{subfigure}{0.48\linewidth}
\centering
\begin{tikzpicture}
\begin{axis}[
    width=\linewidth,
    height=5.2cm,
    domain=-5:5,
    samples=300,
    xlabel={\small $t$},
    ylabel={\small densidad},
    ymin=0, ymax=0.42,
    title={\small Bilateral: $p=0.086$ (dos colas)},
    legend style={draw=none, fill=none, at={(0.97,0.97)}, anchor=north east, font=\footnotesize},
    ticklabel style={/pgf/number format/fixed, font=\footnotesize},
    label style={font=\footnotesize}
]
% Curva (ilustrativa con N(0,1))
\addplot[name path=curve, very thick] {1/sqrt(2*pi)*exp(-x^2/2)};
\def\tcrit{1.72}
% Colas sombreadas (bilateral)
\addplot[fill=blue!30, draw=none, domain=\tcrit:5] {1/sqrt(2*pi)*exp(-x^2/2)} \closedcycle;
\addplot[fill=blue!30, draw=none, domain=-5:-\tcrit] {1/sqrt(2*pi)*exp(-x^2/2)} \closedcycle;
% Líneas verticales
\addplot[dashed] coordinates {(\tcrit,0) (\tcrit,0.41)};
\addplot[dashed] coordinates {(-\tcrit,0) (-\tcrit,0.41)};
\legend{\footnotesize densidad $t$ }
\end{axis}
\end{tikzpicture}
\caption{\footnotesize Prueba bilateral $H_1\neq 0$}
\end{subfigure}\hfill
\begin{subfigure}{0.48\linewidth}
\centering
\begin{tikzpicture}
\begin{axis}[
    width=\linewidth,
    height=5.2cm,
    domain=-5:5,
    samples=300,
    xlabel={\small $t$},
    ylabel={\small densidad},
    ymin=0, ymax=0.42,
    title={\small Unilateral: $p\approx 0.957$ (cola izquierda)},
    legend style={draw=none, fill=none, at={(0.97,0.97)}, anchor=north east, font=\footnotesize},
    ticklabel style={/pgf/number format/fixed, font=\footnotesize},
    label style={font=\footnotesize}
]
% Curva
\addplot[name path=curve2, very thick] {1/sqrt(2*pi)*exp(-x^2/2)};
\def\tcrit{1.72}
% Cola izquierda (H1: beta1 < 0) para t_obs > 0
\addplot[fill=blue!20, draw=none, domain=-5:\tcrit] {1/sqrt(2*pi)*exp(-x^2/2)} \closedcycle;
% Línea vertical en t_obs
\addplot[dashed] coordinates {(\tcrit,0) (\tcrit,0.41)};
\legend{\footnotesize densidad $t$ }
\end{axis}
\end{tikzpicture}
\caption{\footnotesize Prueba unilateral $H_1:\beta_1<0$}
\end{subfigure}

\end{figure}

{\bf Interpretación del valor-$p$ de la cola izquierda.}
Para contrastar $H_0:\beta_1=0$ frente a $H_1:\beta_1<0$, el valor-$p$ de la cola izquierda se define como
\[
\text{valor-p}_{\text{cola izquierda}}=\Pr_{H_0}\!\big(T \le t_{\text{obs}}\big),
\]
es decir, la probabilidad (bajo $H_0$) de observar un estadístico $t$ \emph{tan pequeño como} el observado o menor (solo sirve como evidencia en contra de la nula la cola izquierda).
Si $t_{\text{obs}}>0$, la mayor parte de la masa de la distribución queda a la izquierda de $t_{\text{obs}}$ y, por tanto, es grande.

\[
\text{valor-p}_{\text{cola izquierda}}=1-\text{valor-p}_{\text{cola derecha}}=1-0.043=0.957.
\]
\noindent
Por consiguiente, con $H_1:\beta_1<0$ y $t_{\text{obs}}>0$, el resultado \emph{no} es significativo a niveles habituales (p.\,ej., 10\%).
}

{\bf Ejercicio 2} Considere el siguiente modelo de regresión para el precio de las viviendas:

\[
\log(\text{price}) = \beta_0 + \beta_1 \text{sqrft} + \beta_2 \text{bdrms} + u
\]

\begin{enumerate}
    \item Usted desea estimar y obtener un intervalo de confianza para la \textbf{variación porcentual del precio} (\textit{price}) cuando una casa tiene habitación adicional de 150 pies cuadrados. En forma decimal, esta variación se expresa como:
    \[
    \Delta \log(\text{price}) = 150\beta_1 + \beta_2=\theta_1
    \]

    Sea el vector de coeficientes estimados
\[
\hat{\boldsymbol\beta}=
\begin{pmatrix}
\hat\beta_0\\ \hat\beta_1\\ \hat\beta_2
\end{pmatrix}
=
\begin{pmatrix}
8.10\\ 0.00090\\ 0.050
\end{pmatrix}\]


{con errores estándar } 
\[\text{se}(\hat\beta_0)=0.15,\;
\text{se}(\hat\beta_1)=0.00025,\;
\text{se}(\hat\beta_2)=0.020.
\]

La matriz var--covarianza es
\[
\widehat{\mathrm{Var}}(\hat{\boldsymbol\beta})=
\begin{pmatrix}
0.02250 & -1.00\times 10^{-5} & 3.00\times 10^{-4}\\
-1.00\times 10^{-5} & 6.25\times 10^{-8} & -1.50\times 10^{-6}\\
3.00\times 10^{-4} & -1.50\times 10^{-6} & 4.00\times 10^{-4}
\end{pmatrix}.
\]


Obtenga el \textbf{error estándar} de $\hat{\theta}_1$ y construya  su \textbf{intervalo de confianza del 95\%}.



{\color{blue}

\[
\theta_1 \equiv 150\,\beta_1 + \beta_2, 
\qquad 
\hat\theta_1=150\,\hat\beta_1+\hat\beta_2
=150(0.00090)+0.050=0.185.
\]
Interpretación: incremento aproximado del \(18.5\%\) en \(\textit{price}\).

{Error estándar de \(\hat\theta_1\).}
Con \(\mathbf{a}=(0,\,150,\,1)^\top\), se tiene
\[
\widehat{\mathrm{Var}}(\hat\theta_1)=
\mathbf{a}^\top\,\widehat{\mathrm{Var}}(\hat{\boldsymbol\beta})\,\mathbf{a}
=150^2\,\mathrm{Var}(\hat\beta_1)+\mathrm{Var}(\hat\beta_2)+2\cdot150\,\mathrm{Cov}(\hat\beta_1,\hat\beta_2).
\]
Sustituyendo,
\[
\widehat{\mathrm{Var}}(\hat\theta_1)=
22500\cdot(6.25\times10^{-8})+ (4.00\times10^{-4})
+ 2\cdot150\cdot(-1.50\times10^{-6})
=0.00135625,
\]
\[
\Rightarrow\quad \text{se}(\hat\theta_1)=\sqrt{0.00135625}\approx 0.03683.
\]


Usando cuantiles normales (o \(t\) con g.l. grandes),
\[
\hat\theta_1 \pm 1.96\,\text{se}(\hat\theta_1)
= 0.185 \pm 1.96(0.03683)
= 0.185 \pm 0.0722,
\]
\[
\Rightarrow\quad \text{IC}_{95\%}(\theta_1)\;\approx\;[\,0.113,\;0.257\,].
\]

Con los datos propuestos, la variación porcentual del precio ante una habitación adicional de 150 ft\(^2\) se estima en \(18.5\%\), con \(\text{se}\approx 0.0368\) e intervalo del 95\% \([11.3\%,\,25.7\%]\).
}

\end{enumerate}


{\bf Ejercicio 3}
Considera el modelo lineal \(y=\beta_0+\beta_1 x_1+\beta_2 x_2+u\) con constante y dos regresores. 
Con \(n=5\) observaciones, la matriz de regresores y el vector con la variable dependiente son:
\[
X=\begin{pmatrix}
1&0&1\\
1&1&0\\
1&2&1\\
1&3&0\\
1&4&1
\end{pmatrix},\qquad
Y=\begin{pmatrix}2\\3\\5\\7\\9\end{pmatrix}.
\]

Determine los estimadores de MCO.



{\color{blue} Cálculo de \(X'X\) y \(X'Y\).
A partir de los sumatorios (\(\sum 1=5,\ \sum x_1=10,\ \sum x_2=3,\ \sum x_1^2=30,\ \sum x_2^2=3,\ \sum x_1x_2=6,\ \sum y=26,\ \sum x_1y=70,\ \sum x_2y=16\)), se obtiene:
\[
X'X=
\begin{pmatrix}
5&10&3\\
10&30&6\\
3&6&3
\end{pmatrix},
\qquad
X'Y=
\begin{pmatrix}
26\\ 70\\ 16
\end{pmatrix}.
\]

\paragraph{Inversa de \(X'X\) por cofactores.}
Sea \(A=X'X\). Calculamos su matriz de cofactores \(C=[C_{ij}]\), donde 
\(C_{ij}=(-1)^{i+j}M_{ij}\) y \(M_{ij}\) es el menor (determinante) que resulta de eliminar la fila \(i\) y la columna \(j\) de \(A\).

\subparagraph{Menores y cofactores de la primera fila (\(i=1\)).}
\[
M_{11}=\begin{vmatrix}30&6\\6&3\end{vmatrix}=90-36=54,\ \ \Rightarrow\ C_{11}=+54.
\]
\[
M_{12}=\begin{vmatrix}10&6\\3&3\end{vmatrix}=30-18=12,\ \ \Rightarrow\ C_{12}=-12.
\]
\[
M_{13}=\begin{vmatrix}10&30\\3&6\end{vmatrix}=60-90=-30,\ \Rightarrow\ C_{13}=-30.
\]

\subparagraph{Menores y cofactores de la segunda fila (\(i=2\)).}
\[
M_{21}=\begin{vmatrix}10&3\\6&3\end{vmatrix}=30-18=12,\ \ \Rightarrow\ C_{21}=-12.
\]
\[
M_{22}=\begin{vmatrix}5&3\\3&3\end{vmatrix}=15-9=6,\ \ \Rightarrow\ C_{22}=+6.
\]
\[
M_{23}=\begin{vmatrix}5&10\\3&6\end{vmatrix}=30-30=0,\ \ \Rightarrow\ C_{23}=0.
\]

\subparagraph{Menores y cofactores de la tercera fila (\(i=3\)).}
\[
M_{31}=\begin{vmatrix}10&3\\30&6\end{vmatrix}=60-90=-30,\ \Rightarrow\ C_{31}=-30.
\]
\[
M_{32}=\begin{vmatrix}5&3\\10&6\end{vmatrix}=30-30=0,\ \Rightarrow\ C_{32}=0.
\]
\[
M_{33}=\begin{vmatrix}5&10\\10&30\end{vmatrix}=150-100=50,\ \Rightarrow\ C_{33}=+50.
\]

\subparagraph{Matriz de cofactores, adjunta y determinante.}
\[
C=\begin{pmatrix}
54 & -12 & -30\\
-12& \ \ 6 & \ \ 0\\
-30& \ \ 0 & 50
\end{pmatrix},
\qquad
\operatorname{adj}(A)=C^\top=\begin{pmatrix}
54 & -12 & -30\\
-12& \ \ 6 & \ \ 0\\
-30& \ \ 0 & 50
\end{pmatrix}.
\]
El determinante (por expansión en la primera fila) es:
\[
\det(A)=5\cdot 54+10\cdot(-12)+3\cdot(-30)=270-120-90=60\neq 0.
\]

\subparagraph{Inversa.}
\[
A^{-1}=\frac{\operatorname{adj}(A)}{\det(A)}
=\frac{1}{60}\begin{pmatrix}
54 & -12 & -30\\
-12& \ \ 6 & \ \ 0\\
-30& \ \ 0 & 50
\end{pmatrix}
=
\begin{pmatrix}
0.9 & -0.2 & -0.5\\
-0.2& 0.1 & 0\\
-0.5& 0 & 5/6
\end{pmatrix}.
\]

\paragraph{Estimador MCO.}
\[
\hat\beta=(X'X)^{-1}X'Y
=
\begin{pmatrix}
0.9 & -0.2 & -0.5\\
-0.2& 0.1 & 0\\
-0.5& 0 & 5/6
\end{pmatrix}
\begin{pmatrix}
26\\ 70\\ 16
\end{pmatrix}
=
\begin{pmatrix}
1.4\\ 1.8\\ 1/3
\end{pmatrix}.
\]

\paragraph{Ecuación estimada.}
\[
\hat y \;=\; 1.4 \;+\; 1.8\,x_1 \;+\; \tfrac{1}{3}\,x_2.
\]

}

\end{document}
